\documentclass[11pt,a4paper]{article}
\usepackage[margin=27mm]{geometry}
\usepackage{fontspec,unicode-math}
\setmainfont{STIXGeneral}
\setmathfont{STIX Two Math}
\usepackage{amsmath,amsthm,booktabs,hyperref}
\hypersetup{hidelinks}
\newtheorem{theorem}{Theorem}
\newtheorem{lemma}[theorem]{Lemma}
\title{A bipartite graph with a nonsplit action on its parts}
\author{Finite construction and elementary certificate}
\date{12 October 2026}
\begin{document}
\maketitle
\begin{abstract}
We give an explicit connected cubic bipartite graph whose automorphisms
preserving its parts act transitively on edges, but whose automorphism group
does not split over that subgroup.  The graph and two permutations are
specified by complete finite data.  An elementary certificate proves that
these permutations generate the full automorphism group, so the absence of
part-exchanging involutions is not merely a property of a chosen subgroup.
\end{abstract}

\section{The construction and its finite data}
Let $B$ be the graph on $\{1,\ldots,896\}$ with adjacency lists in
\texttt{graph.json}.  A row lists exactly the three neighbours of the vertex
with that row number.  The same file gives two permutations by their complete
image lists.  All indices in this data file start at one.  These lists, rather
than a random search or a chosen graph-library labelling, define the graph.

Two whole files form the required certificate: \texttt{check.py} and
\texttt{graph.json}.  Only Python's standard library is used.  The successful
report, \texttt{certificate.json}, binds both files by SHA256.  The checks use
explicit conditional failures and remain active under Python's \texttt{-O}
option.  The certificate does not call GAP, nauty, Bliss, or a graph
isomorphism package.

\begin{theorem}\label{thm:graph}
The graph $B$ is connected, simple, cubic and bipartite, with parts $P,Q$ of
size $448$ and $1344$ edges.  Write $A=\operatorname{Aut}(B)$ and let $A^+$
be the subgroup preserving each part.  Then
\[
 |A|=5376,\qquad |A^+|=2688,
\]
and $A^+$ is transitive on the edges.  The elements of $A\setminus A^+$
have the following orders and multiplicities:
\[
\begin{array}{c|rrr}
\text{order}&4&12&28\\ \hline
\text{multiplicity}&1024&896&768.
\end{array}
\]
In particular, the sequence
$1\longrightarrow A^+\longrightarrow A\longrightarrow C_2\longrightarrow1$
does not split.
\end{theorem}

The proof below separates the small general argument certifying fullness
from the finite calculations.  The latter are required premises and their
whole data and implementation are retained.  Optional independent computations
with nauty and Bliss agree, but are not needed for the elementary proof.

\section{Why the computed group is the full group}
For a finite graph and an initial assignment of integer colours, replace the
colour of each vertex by the pair consisting of its previous colour and the
multiset of the previous colours of its neighbours.  Encode equal pairs by
equal new integers, using a fixed sorted order.  This is neighbour-colour
refinement.  It never merges two previous colour classes.

\begin{lemma}\label{lem:bound}
Let $X$ be a finite $d$-regular graph on $n$ vertices.  Suppose that giving
distinct colours to one vertex $v$ and its $d$ neighbours, and one additional
colour to all other vertices, eventually produces singleton classes under
neighbour-colour refinement.  Then $|\operatorname{Aut}(X)|\leq n d!$.
\end{lemma}
\begin{proof}
An automorphism fixing $v$ and every neighbour of $v$ preserves the initial
colours.  Inductively it preserves each refined colour: it preserves both
the old colour and the multiset of old colours in the neighbourhood.  When
all colour classes are singletons, the automorphism fixes every vertex.
Consequently an automorphism is determined by its action on the ordered
star consisting of $v$ and its neighbours.  There are at most $n$ choices
for the image of $v$ and $d!$ bijections from its neighbourhood to that of
the image.  This proves the bound.
\end{proof}

\begin{proof}[Proof of Theorem~\ref{thm:graph}]
The certificate first checks that every adjacency row contains three
distinct valid vertices, excludes its own vertex, and is reciprocal.
Breadth-first traversal visits all $896$ vertices, assigns opposite binary
colours across every edge and finds $448$ vertices of each colour.  Counting
unordered adjacent pairs gives $1344$ edges.

Individualize vertex $1$ and its three neighbours.  Starting with five
colour classes, successive strict refinement steps give
\[
5,8,11,14,17,20,26,44,135,398,692,842,890,896.
\]
The final partition is discrete.  Lemma~\ref{lem:bound} therefore gives
$|A|\leq896\cdot6=5376$.

The two supplied permutations are bijections of the vertex set and preserve
every adjacency row.  Their orders are $6$ and $4$.  Starting at the identity,
the certificate closes a set of permutations under multiplication by these
two generators and obtains exactly $5376$ distinct permutations.  This
finite closed monoid of permutations is their generated group: every inverse
is a positive power in the finite permutation group.  It is contained in
$A$ and attains the upper bound, so it equals $A$.

For every one of these $5376$ permutations, the certificate tests its image
of $P$.  Exactly $2688$ preserve $P$, and the other $2688$ send it to $Q$.
The orbit of one unordered edge under the former permutations consists of
all $1344$ edges.  For each of the remaining permutations its order is
computed as the least common multiple of the lengths of its disjoint cycles.
The resulting distribution is the one displayed in the theorem; its entries
sum to $2688$.

The part action is onto $C_2$, with kernel $A^+$.  A section of this action
would send the nonidentity element of $C_2$ to an involution outside $A^+$.
The exhaustive distribution contains no such element.  Conversely such an
involution would generate a complement, so the stated obstruction is exact.
\end{proof}

\newpage
\section{Scope and reproducibility}
The graph supplies an input for a separate construction of a thin
flag-transitive incidence geometry.  This note proves the graph assertions
only; it does not assert a source-wide result about incidence geometries or
replace the separate proof identifying the full type-permuting group of
that construction.  Its distinction between edge transitivity under $A^+$
and under $A$ is essential.

The graph was found as an eight-fold binary homology cover of a cubic graph
on $112$ vertices.  That search is provenance, not a premise: the retained
adjacency and permutations are sufficient to replay the entire proof.
In particular, no completeness claim about the search is made.  General
permutation-group arithmetic, breadth-first traversal and neighbour-colour
refinement are standard procedures; no historical novelty claim is made
for these methods.

Seven controls reject a loop, a nonbijective proposed generator, a
permutation not preserving adjacency, substitution of a proper subgroup,
an incomplete star, the genuine swapping involutions of a six-cycle, and
the weakened use of root-only colour refinement.  The six-cycle control
passes all preceding full-group and transitivity tests, so it specifically
tests the final nonsplitting obstruction.

To replay the required certificate from an exported copy, run
\begin{quote}\ttfamily\small
python3 -B -O check.py --out fresh-result.json
\end{quote}
There is no Lean formalization in this packet.  The finite certificate is
an ordinary executable proof by exhaustive verification; both its complete
program and its mandatory data must be counted wherever this construction
is used.  Optional GAP/nauty/Bliss checks are separately inventoried as
validation and cannot stand in for the required fullness argument.
\end{document}
